\section{Crystalline realization and the section problem (Part J)}\label{sec:categoryJ}

\subsection{Data, conventions, and readings decided on the scans}
We use $K$ for a complete discretely valued field of characteristic zero
with perfect residue field $k$ of characteristic $p>0$,
$L=K_0=W(k)[1/p]$, and assume $K/K_0$ finite and totally ramified.
Put $\pi=\operatorname{Gal}(\overline K/K)$,
$\pi_0=\operatorname{Gal}(\overline K/K_0)$,
and $C=\widehat{\overline K}$.  We normalize $v_p(p)=1$.
Scan 023 explicitly says that $k$ is perfect; it does not assume that
$k$ is algebraically closed.

The supplied archive contains the Pro transcription of 161-5/023--029,
but not the second transcription \texttt{K161-5-02} cited in the study.
We subsequently retrieved the original Montpellier PDF and inspected
scans 023--029 directly.  Scan 023 does define
$\mathcal M_K^{(1)}$ as all Hodge--Tate representations of weights $0,1$,
includes actual Barsotti--Tate representations as a full subcategory,
and then defines $\mathcal M_K^{\mathrm{BT}}$ using
$\mathcal M_K^{(1)}$ as generators.  Thus the larger reading is in the
manuscript, not merely in the transcription.  We distinguish it from
the category generated by actual Barsotti--Tate representations.
Theorem~\ref{cat:J7} proves uniqueness on the latter category;
Theorem~\ref{cat:J8} refutes existence on the literal larger category.

There is also an unannounced change of field letters.  The following
translation makes all the later torsor and Frobenius data type-check:
\[
\begin{array}{c|cc}
 &\text{ramified field}&\text{unramified field}\\ \hline
 \text{scan 023}&K&L\\
 \text{scans 024--029, consistent interpretation}&L&K\\
 \text{our notation throughout}&K&K_0=L.
\end{array}
\]
The scan-024 expression $T_{\rm cr}\otimes_K L$ and scan-027
``descent from $L$ to $K$'' are both correctly typed in the second
row.  If one retains scan-023's letters instead, they must respectively
be corrected to $\otimes_LK$ and descent from $K$ to $L$.
The scans do not decide whether these are slips or an unannounced
change of notation.  We do not silently mix the two conventions.

Let $\mathcal C$ be a neutral $\mathbf Q_p$-linear Tannakian category
of Hodge--Tate representations, let $\omega_p$ be its forgetful functor,
and put $G=\operatorname{Aut}^{\otimes}(\omega_p)$.
Evaluation gives the continuous homomorphism
$\rho:\pi\to G(\mathbf Q_p)$.
The Hodge--Tate fibre functor
$\omega_{\rm HT}:\mathcal C\to\operatorname{GrVect}_K$ gives the
right $G_K$-torsor
\[
 Q_K=\operatorname{Isom}^{\otimes}
       (K\otimes_{\mathbf Q_p}\omega_p,\omega_{\rm HT}^{\rm ungr}),
 \qquad G''_K=\operatorname{Aut}_{G_K}(Q_K).
\]
The grading is exactly a cocharacter
$j_K:\mathbf G_{m,K}\to G''_K$.
This is the later manuscript's $Q_L,j_L$, after the field dictionary.
Its earlier description starts with the tensor-isomorphism bitorsor
$\Phi_C$ for the scalar and diagonal $\pi$-actions on
$C\otimes\omega_p$, a cocharacter
$j^!:\mathbf G_{m,C}\to\operatorname{Aut}_{G_C}(\Phi_C)$,
and the torsor $\tau_C$ of Tate-line frames.  With the dual Tate-line
convention that cancels the Hodge--Tate weights, the contracted product
$\tau_C\times^{\mathbf G_{m,C}}\Phi_C$ is $Q_C$.
The assertion called an axiom on scan 026 is precisely the required
field-of-definition descent of this graded torsor; in the Hodge--Tate
realization above the descended object is $Q_K$.

Let $U''_K\subset G''_K$ be the unipotent subgroup acting trivially
on the graded of the filtration induced by $j_K$.  A right
$U''_K$-torsor $R_K$ defines the right $G_K$-torsor
\[
 P_K=R_K\times^{U''_K}Q_K.
\]
The contracted product is well defined because $U''_K$ acts on $Q_K$
by $G_K$-equivariant automorphisms.  It describes the filtered fibre
functor with its identification of the associated graded with
$\omega_{\rm HT}$.  These are the later manuscript's $R_L,P_L$.
The next datum is descent of the underlying torsor from $K$ to $K_0$:
a right $G_{K_0}$-torsor $P_{K_0}$ and an isomorphism
$P_{K_0}\times_{K_0}K\simeq P_K$.
The filtration remains over $K$.  Write
$G'_{K_0}=\operatorname{Aut}_{G_{K_0}}(P_{K_0})$.
Finally, with $\sigma$ Witt Frobenius on $K_0$, Frobenius is an
isomorphism $\sigma^*P_{K_0}\simeq P_{K_0}$.
Thus the later manuscript's $P_K\simeq P_K^{(\varphi)}$ is the
unramified Frobenius datum, not a Frobenius on the filtered field $K$.

A concrete group encoding field descent and Frobenius is
\[
 \Gamma=\{(n,\gamma):n\in\mathbf Z,
   \ \gamma\in\operatorname{Aut}_{\rm cont}(C),
   \ \gamma|_{K_0}=\sigma^n\}.
\]
Composition is $(n,\gamma)(m,\delta)=(n+m,\gamma\delta)$.
An automorphism of $K_0$ extends to the algebraic closure and its
completion, so
$1\to\pi_0\to\Gamma\to\mathbf Z\to1$ is exact.
The subgroup whose already specified action must be retained is
$\pi\subset\pi_0$, not necessarily the entire kernel.
In particular $\pi$ need not be normal in $\Gamma$; when $K/K_0$
is not Galois, it is not even normal in $\pi_0$.
For $S=\operatorname{Spec}C$ and $P_C$, the group of semilinear
$G_C$-equivariant lifts gives the extension
\[
 1\longrightarrow G'_C(C)\longrightarrow E\longrightarrow\Gamma
 \longrightarrow1
\]
whenever these lifts exist.  The given descent action is its fixed
partial section $\alpha_0:\pi\to E$.
Its invariant automorphism group is
$G'_C(C)^\pi=G'_K(K)$, which is exactly the later manuscript's
$G'_L(L)$.  Equality follows by descent (or on coefficients by the
fixed-field theorem and the affine defining equations).
Only in cases such as $K=K_0$ is the quotient by this fixed subgroup
$\pi$ the infinite cyclic group and its invariant field $K_0$.
The relative classification below deliberately does not assume
normality.

\subsection{The complete abstract classification}
Let $S$ be a ring with an action of a group $\Gamma$, let $G_S$ be a group
scheme with compatible semilinear $\Gamma$-action, and let $P$ be a right
$G_S$-torsor.  Assume the semilinear lifts exist for each element of
$\Gamma$.  The group $E$ of such lifts fits into
\[
 1\longrightarrow H\longrightarrow E\longrightarrow\Gamma
 \longrightarrow1,\qquad H=\operatorname{Aut}_{G_S}(P)(S).
\]
This is the extension of cote 161-5/028.  For topological groups all sections
and cocycles below are required to be continuous; the same proofs apply.

\begin{theorem}[Proved: exact classification of extensions]\label{cat:J1}
Fix a subgroup $\pi\subset\Gamma$ and a section
$\alpha_0:\pi\to E|_\pi$.  If there is no extending section, the solution
set is empty.  If an extending section $\alpha:\Gamma\to E$ exists, let
$\gamma\cdot h=\alpha(\gamma)h\alpha(\gamma)^{-1}$ be the resulting action
on $H$.  Then the groupoid of extending sections, with conjugations that
preserve $\alpha_0$, is
\[
 \bigl[Z^1(\Gamma,H;\pi)/H^\pi\bigr],
 \quad
 Z^1(\Gamma,H;\pi)=
 \{c:\Gamma\to H:\ c(\gamma\delta)=c(\gamma)
       (\gamma\cdot c(\delta)),\ c|_\pi=1\}.
\]
Here $h\in H^\pi$ acts by
$c(\gamma)\mapsto h c(\gamma)(\gamma\cdot h^{-1})$.
If $\pi$ is normal, this groupoid is canonically
$[Z^1(\Gamma/\pi,H^\pi)/H^\pi]$; its set of isomorphism classes is
$H^1(\Gamma/\pi,H^\pi)$.  The identifications depend on the reference
section $\alpha$; changing it twists the coefficient action.
\end{theorem}
\begin{proof}
For another section $\beta$ put
$c(\gamma)=\beta(\gamma)\alpha(\gamma)^{-1}\in H$.
Multiplying $\beta(\gamma)\beta(\delta)$ gives the displayed cocycle
identity.  Conversely that identity says precisely that
$\beta(\gamma)=c(\gamma)\alpha(\gamma)$ is a homomorphism section.
Equality with $\alpha_0$ on $\pi$ is precisely $c|_\pi=1$.
Conjugating $\beta$ by $h$ gives the stated formula.  Such a conjugation
preserves $\alpha_0$ if and only if $h$ commutes with $\alpha_0(\pi)$,
which is the condition $h\in H^\pi$.  Thus even a conjugation allowed
initially in all of $H$ lies in $H^\pi$ when its source and target both
extend the fixed $\alpha_0$.

Assume $\pi$ is normal.  For $n\in\pi$, the cocycle identity gives
$c(\gamma n)=c(\gamma)$ and
$c(n\gamma)=n\cdot c(\gamma)$.  Since
$n\gamma=\gamma(\gamma^{-1}n\gamma)$, the first identity shows that
$n\cdot c(\gamma)=c(\gamma)$.  Consequently $c$ takes values in $H^\pi$
and is constant on cosets.  The $\Gamma$-action on $H^\pi$ factors
through $\Gamma/\pi$, so $c$ descends to a cocycle on the quotient.
Inflation is inverse to this construction and preserves the action of
$H^\pi$.  For continuous cocycles the descended map is continuous by
the quotient topology.  This proves all assertions, including morphisms
and automorphisms in the groupoids.
\end{proof}

\begin{corollary}[Proved: Frobenius classes]\label{cat:J2}
Suppose $\pi\triangleleft\Gamma$, $\Gamma/\pi=\mathbf Z$, and
$H^\pi=G'_L(L)$.  Write $\sigma$ for the action of $1\in\mathbf Z$ on
$G'_L(L)$ induced by a fixed extending section.  The extensions are
classified by the $\sigma$-conjugacy classes
\[
 G'_L(L)/\bigl(b\sim h b\sigma(h)^{-1}\bigr).
\]
The automorphism group of the class represented by $b$ is
$\{h\in G'_L(L):h b=b\sigma(h)\}$.
\end{corollary}
\begin{proof}
A cocycle on the infinite cyclic group is uniquely determined by its
value $b$ at $1$; at $n>0$ it is
$b\sigma(b)\cdots\sigma^{n-1}(b)$, and its values at negative integers
are forced by the cocycle identity.  Every $b$ gives such a cocycle.
The equivalence and stabilizer formulas are exactly those in
Theorem~\ref{cat:J1}.
\end{proof}

\begin{proposition}[Refuted: slopes or weak admissibility do not imply uniqueness]\label{cat:J3}
The lead H1615-1 is false under the stated perfect-residue-field
hypothesis.  For $k=\mathbf F_p$, $L=K=\mathbf Q_p$, $G'=\mathbf G_m$,
and Witt Frobenius $\sigma=1$, the classes of
Corollary~\ref{cat:J2} are $\mathbf Q_p^\times$, whereas the Newton slope
is only $v_p:\mathbf Q_p^\times\to\mathbf Z$.
For every $a\in\mathbf Z_p^\times$, the rank-one filtered isocrystal
\[
 D_a=\mathbf Q_p e,\qquad \varphi(e)=a e,\qquad
 \operatorname{Fil}^iD_a=
 \begin{cases}D_a&i\leq0,\\0&i>0\end{cases}
\]
has slope zero and is weakly admissible.  Two such objects are
isomorphic if and only if their parameters are equal.  In particular
$a=1$ and $a=1+p$ give distinct classes with identical slopes and
identical filtration.
\end{proposition}
\begin{proof}
Every $\sigma$-coboundary is $h/\sigma(h)=1$, so the cyclic cohomology
quotient is the whole multiplicative group.  This quotient occurs with exactly the fixed data above when
$K=K_0=\mathbf Q_p$: take $S=C$, $\pi=\operatorname{Gal}(\overline K/K)$,
$\Gamma=\pi\times\mathbf Z$, and let the second factor act trivially
on $C$.  This is the displayed definition of $\Gamma$, since
$\sigma=1$ on $K_0$.  One may take $G=\mathbf G_m$ and trivial $\rho$, as the abstract
section problem permits.  To require also that $G$ be the actual
Tannaka group of a representation, take instead the unramified
character $\rho$ sending arithmetic Frobenius to $1+p$.
Its image is infinite and hence Zariski dense in $\mathbf G_m$.
It is the Tate representation of an etale $p$-divisible group: at
each level descend the constant cyclic group using the unramified
action on $\mathbf Z/p^n\mathbf Z$, and take the compatible system.
Thus it has Hodge--Tate weight zero.
In either version $j=0$, $U''=1$, $R_K$ is the point, and $Q_K=P_K$
is the rank-one Hodge--Tate frame torsor, which is neutral over $K$.
In a $C$-frame with its original descent action,
$E=C^\times\rtimes\Gamma$ and
$\alpha_0(g)=(\rho(g),(g,0))$.
For each $a\in\mathbf Q_p^\times$ the section
$(g,n)\mapsto(a^n\rho(g),(g,n))$ extends this same $\alpha_0$ and
retains all the same $\rho,j,Q,R,P$ data.
The fixed-field theorem gives $H^\pi=\mathbf Q_p^\times$.
Continuity holds with the discrete topology on $\mathbf Z$.
Thus the counterexample is within the manuscript's field and section
scope, and not merely a change of the fixed partial section.


For $D_a$, $t_N(D_a)=v_p(a)=0=t_H(D_a)$.  Its only proper subobject
is zero, so all weak-admissibility inequalities hold.  An isomorphism
$D_a\to D_b$ must send $e$ to $c e$ for some $c\in\mathbf Q_p^\times$.
Commutation with Frobenius requires $ac=bc$, whence $a=b$.
Thus neither equal slopes nor weak admissibility removes the ambiguity.
This is also a counterexample to conjugacy under the full group $H$,
by the last assertion of the first paragraph of
Theorem~\ref{cat:J1}'s proof.
\end{proof}

\begin{proposition}[Proved: the rank-one classification over a perfect residue field]\label{cat:J4}
Let $k$ be any perfect field of characteristic $p$, $L=W(k)[1/p]$, and
$\sigma$ Witt Frobenius.  Rank-one isocrystals over $L$ are classified by
\[
 L^\times/\{\sigma(c)/c:c\in L^\times\}.
\]
There is a split exact sequence of abelian groups
\[
 1\longrightarrow W(k)^\times/\{\sigma(u)/u:u\in W(k)^\times\}
 \longrightarrow H^1(\mathbf Z,L^\times)
 \xrightarrow{v_p}\mathbf Z\longrightarrow0,
\]
whose splitting sends $n$ to $[p^n]$.  If $k$ is algebraically closed,
the left term is trivial, so slopes classify rank-one isocrystals.
\end{proposition}
\begin{proof}
In a basis, Frobenius is $e\mapsto a e$ semilinearly, and replacing
the basis by $c e$ replaces $a$ by $a\sigma(c)/c$.  Each nonzero
$c\in L$ is uniquely $p^n u$ with $u\in W(k)^\times$, and
$\sigma(p)=p$.  The decomposition $L^\times=p^{\mathbf Z}\times
W(k)^\times$ proves exactness and splitting.

Suppose $k$ is algebraically closed and let $a\in W(k)^\times$.
Choose $c_1\in W(k)^\times$ such that
$\sigma(c_1)/c_1\equiv a\pmod p$; modulo $p$ this asks for a
nonzero root of $X^{p-1}=\overline a$.
Inductively assume the congruence holds modulo $p^n$, $n\geq1$.
Writing $c_{n+1}=c_n(1+p^n[x])$, the ratio changes modulo $p^{n+1}$
by the factor $1+p^n[x^p-x]$.
The required correction therefore reduces to one equation
$x^p-x=b$ in $k$, which has a root.  Choose such a correction at
each step.  The units $c_n$ form a $p$-adically Cauchy sequence;
its limit $c$ satisfies $\sigma(c)/c=a$.  This proves surjectivity
on units and the last assertion.  The argument works also for $p=2$.
\end{proof}

\subsection{The positive realization statement and its precise uniqueness}
\begin{theorem}[Known, with the hypotheses checked here]\label{cat:J5}
Let $\mathcal C_{\mathrm{BT}}$ be the smallest full Tannakian subcategory
of $\operatorname{Rep}_{\mathbf Q_p}(\pi_K)$ containing the rational
Tate modules of $p$-divisible groups over $\mathcal O_K$.
Then
\[
 D_{\mathrm{cris}}(V)=
 (B_{\mathrm{cris}}\otimes_{\mathbf Q_p}V)^{\pi_K}
\]
is an exact tensor functor on $\mathcal C_{\mathrm{BT}}$ to admissible
filtered $\varphi$-modules over $(L,K)$.  On Barsotti--Tate objects it
is their crystalline realization, with the covariant or contravariant
Dieudonne convention matched by dualizing as necessary.  Its de Rham
realization is $K\otimes_LD_{\mathrm{cris}}(V)$, and the associated
graded is the Hodge--Tate realization.
\end{theorem}
\begin{proof}[Derivation from opened theorems]
Fontaine~\cite[\S6.2]{catFontaine1982} proves that rational Tate
modules of $p$-divisible groups over $\mathcal O_K$ are crystalline
and identifies the filtered Dieudonne realization with the period
realization.  His theorem in \S5.2(i) proves that crystalline
representations are closed under the operations defining
$\mathcal C_{\mathrm{BT}}$; \S5.2(iii)--(v) gives the exact tensor
equivalence and the de Rham scalar-extension comparison.
The associated-graded comparison is the Hodge--Tate comparison of
\S3, in particular \S3.6--3.7 and \S3.10.  Thus all these assertions restrict to
$\mathcal C_{\mathrm{BT}}$.  We have used the actual BT-generated
subcategory; no assertion about all Hodge--Tate representations is
needed in this deduction.
\end{proof}

\begin{proposition}[Proved: uniqueness with the crystalline comparison]\label{cat:J6}
On $\mathcal C_{\mathrm{BT}}$, consider tensor functors $F$ to
admissible filtered $\varphi$-modules together with a specified tensor
natural isomorphism
$\eta:V_{\mathrm{cris}}\circ F\xrightarrow{\sim}
\operatorname{incl}_{\mathcal C_{\mathrm{BT}}}$.
The groupoid of such pairs is contractible: there is an isomorphism
to $(D_{\mathrm{cris}},\text{canonical comparison})$, unique among
isomorphisms compatible with $\eta$.  Weak admissibility may replace
admissibility in this statement.
\end{proposition}
\begin{proof}
Fontaine's equivalence in \S5.2 identifies
$D_{\mathrm{cris}}V_{\mathrm{cris}}F$ with $F$ by its counit.
Applying $D_{\mathrm{cris}}$ to $\eta$ and composing with the inverse
of this counit produces $F\simeq D_{\mathrm{cris}}|_{\mathcal C_{\mathrm{BT}}}$.
The construction is tensor and natural because the equivalence, counit,
and $\eta$ are.  Faithfulness of $V_{\mathrm{cris}}$ shows that at most
one such isomorphism induces the specified $\eta$.
Colmez--Fontaine~\cite[Theorem A; \S4.1, Theorem 4.3 and its
corollary]{catColmezFontaine2000} identifies weakly admissible
filtered $(\varphi,N)$-modules with admissible ones; setting $N=0$
gives the crystalline assertion.  This replacement preserves the
comparison datum, rather than deducing it from the Newton slopes.
\end{proof}

The field hypotheses of these sources match our scope.  Fontaine's
Introduction, printed p.~529, assumes a perfect residue field, not a
finite one.  Colmez--Fontaine's Introduction, p.~2 of the opened author
version, assumes a perfect field $k$ and a finite totally ramified
extension of $W(k)[1/p]$.  Kisin makes the same assumptions at the
start of his Introduction.  Thus none of the three applications here
silently narrows the manuscript to finite extensions of $\mathbf Q_p$.
These are statements at the rational or isogeny level; no assertion
about a chosen integral lattice at $p=2$ is being substituted for them.

Kisin~\cite[Theorem (0.3), Corollary (2.2.6)]{catKisin2006} proves
that every crystalline representation of Hodge--Tate weights $0,1$
is the rational Tate module of a $p$-divisible group, with all primes
included at the isogeny level.  This identifies the weight-$[0,1]$
part of the crystalline category with the BT objects.  It does not
identify all Hodge--Tate representations of these weights with them.


\begin{theorem}[Proved: uniqueness from the listed graded comparison]\label{cat:J7}
On $\mathcal C_{\mathrm{BT}}$, let $F$ be an exact tensor functor to
finite filtered $\varphi$-modules over $(L,K)$, with a specified tensor
natural isomorphism
\[
 \gamma:\operatorname{gr}(K\otimes_LF(-))\xrightarrow{\sim}D_{\mathrm{HT}}(-).
\]
Suppose its restriction to actual BT representations is identified with
the BT crystalline functor, compatibly with $\gamma$.  Then there is
exactly one tensor natural isomorphism
$F\simeq D_{\mathrm{cris}}|_{\mathcal C_{\mathrm{BT}}}$ compatible
with these data.  In particular all its values are admissible; that
condition need not be imposed in advance on arbitrary objects.
\end{theorem}
\begin{proof}
We first record the rigidity used in the proof.  On admissible filtered
$\varphi$-modules the associated-graded functor is faithful.  Indeed,
under Fontaine's equivalence its composite is $D_{\mathrm{HT}}$ on
crystalline representations, which is faithful by
\cite[Proposition 1.6(iii)]{catFontaine1982}.  Thus a morphism between
two admissible objects is determined by its associated graded.

Call a finite direct sum of tensor products of BT objects and their
duals a word.  The given comparison on BT objects produces a comparison
$\eta_W:F(W)\simeq D_{\mathrm{cris}}(W)$ on each word.  Both sides are
admissible, since tensor products, duals and sums of crystalline objects
are crystalline.  For every morphism between words, the two composites
with these comparisons have the same associated graded, by naturality
of $\gamma$.  Faithfulness just proved makes the composites equal.
The same holds for direct summands of words, by applying this equality
to the defining idempotents.  In particular it holds for their exterior
powers.

Every rank-one object $T$ of $\mathcal C_{\mathrm{BT}}$ is crystalline
and has a single integral Hodge--Tate weight.  Tensoring it with a
suitable integral power of the Tate object makes that weight zero.
Kisin's Corollary (2.2.6) then makes the twist an actual BT
representation.  The Tate object or its dual is itself BT.
Consequently $T$ is isomorphic to a word.  This gives a comparison
$\eta_T$ with identity associated graded; the same faithfulness makes
it independent of the chosen expression and compatible with all maps
between such rank-one objects and direct summands of words.

Let $X\subset W$ be any subobject of a word and put $r=\dim X$.
For $r>0$, its determinant inclusion
\[
 \det X\longrightarrow\Lambda^r W
\]
is a morphism between the objects for which the comparison is already
known.  The comparison therefore identifies the two determinant lines
in $\Lambda^rF(W)$ and $\Lambda^rD_{\mathrm{cris}}(W)$.
A rank-$r$ subspace $U$ of a vector space is recovered from its
nonzero determinant vector $u$ by
$U=\{w:w\wedge u=0\}$; this follows by completing any basis of $U$
to a basis of the ambient space and reading the remaining coefficients.
Thus $\eta_W$ identifies the images of $F(X)$ and
$D_{\mathrm{cris}}(X)$.  The case $r=0$ is immediate.
Their filtrations are the induced filtrations because the functors
are exact (equivalently, here strictness follows from the exact natural
graded comparison).  Restriction gives a filtered $\varphi$-isomorphism
$\eta_X$, with the prescribed associated graded.

Every object $V$ of the generated category is a subquotient of a word.
Write it as $X/Y$ with $Y\subset X\subset W$ and use the preceding
comparisons and exactness to obtain $\eta_V$ on the quotient.
This also proves that $F(V)$ is admissible.  All comparisons so obtained
have identity associated graded under $\gamma$; faithfulness on
admissible objects now proves that they are independent of the chosen
presentation, natural for every morphism, compatible with tensor
products, and unique.  This proves every assertion.
\end{proof}

\begin{theorem}[Refuted: the literal all-Hodge--Tate extension]\label{cat:J8}
The extension asked for in 161-5/024 does not exist in general if
$\mathcal M_K^{\mathrm{BT}}$ is interpreted literally as the category
generated by all Hodge--Tate representations of weights $0,1$, as on
scan 023.  Already for $K=L=\mathbf Q_5$, no tensor functor with its
specified natural graded Hodge--Tate comparison can agree with the
BT crystalline functor on all actual BT objects and be defined on
that larger category.  This contradiction does not assume weak
admissibility of its other values.
\end{theorem}
\begin{proof}
Let $E=\mathbf Q_5(5^{1/3},\zeta_3)$.  The polynomial $X^3-5$ is
Eisenstein, so adjoining $5^{1/3}$ gives a totally ramified extension
of degree three.  The polynomial $X^2+X+1$ is irreducible modulo five
and separable, so adjoining $\zeta_3$ gives the unramified quadratic
extension.  These extensions have trivial intersection.  Their
composite is the splitting field of $X^3-5$ and has degree six.
The automorphisms
\[
 r(5^{1/3})=\zeta_3 5^{1/3},\quad r(\zeta_3)=\zeta_3,
 \qquad
 s(5^{1/3})=5^{1/3},\quad s(\zeta_3)=\zeta_3^{-1}
\]
satisfy $r^3=s^2=1$ and $srs=r^{-1}$.  Hence
$G=\operatorname{Gal}(E/K)=S_3$, with inertia $\langle r\rangle$ and
unramified quotient $G/\langle r\rangle=C_2$.

Every representation in $\operatorname{Rep}_{\mathbf Q_5}(G)$ is
Hodge--Tate of weight zero, and its Hodge--Tate fibre functor is
\[
 \omega_E(V)=(E\otimes_KV)^G.
\]
Indeed it becomes trivial over $E$; Tate's fixed-field theorem gives
$C^{\operatorname{Gal}(\overline K/E)}=E$, and finite Galois descent
gives the displayed formula and its full dimension
\cite[main Theorem, p.~417; proof in \S3]{catAx1970}.
It is therefore a full subcategory of the literal category in question.

For these objects the graded comparison forces the filtration of $F$
to be concentrated in degree zero and identifies its underlying fibre
functor with $\omega_E$.  Since Witt Frobenius on $L=\mathbf Q_5$ is
the identity, the Frobenius maps of $F$ form a tensor natural
automorphism of $\omega_E$.
Such an automorphism must be the identity.  To see this without any
unstated Tannakian identification, take the regular function algebra
$K^G$ with its left regular $G$-action.  It is a commutative algebra
object of $\operatorname{Rep}_K(G)$, and its descended algebra is
$\omega_E(K^G)\simeq E$.  A tensor automorphism induces a $K$-algebra
automorphism of $E$.  Naturality makes it commute with every right
translation of $K^G$, whose descents give the full group
$\operatorname{Gal}(E/K)$.  It therefore lies in $Z(S_3)=1$.
Every finite-dimensional $G$-representation embeds into a finite direct
sum of regular representations (average a splitting, or use the usual
matrix-coefficient embedding).  Naturality then makes the tensor
automorphism the identity on every object.

Now let $\varepsilon$ be the sign character of $S_3$.  It factors
through the unramified quadratic quotient.  It is an actual BT
representation: descend the constant etale $5$-divisible group
$\mathbf Q_5/\mathbf Z_5$ from the unramified quadratic extension by
the involution $-1$ at each finite level.  Finite etale Galois descent
is effective, and the compatible level maps give an etale
$5$-divisible group over $\mathbf Z_5$.
Its crystalline Frobenius is $-1$.  For an explicit computation put
$u=\zeta_3-\zeta_3^{-1}$ in the unramified quadratic field.  Witt
Frobenius sends $u$ to $-u$, and the invariant generator
$u\otimes e_\varepsilon$ of the corresponding rank-one crystalline
module consequently has Frobenius $-1$; the filtration has degree zero.
This is also the etale BT comparison in
\cite[\S6.2]{catFontaine1982}.
The proposed functor would give this object Frobenius both $1$ and
$-1$, a contradiction.

The exact program \texttt{code/category\_s3\_obstruction.py} enumerates
all six permutations, checks all $36$ products and $216$ associativity
triples, verifies the center and sign action, and checks the stated
residue-polynomial and Eisenstein conditions.  Its saved output agrees
with each finite part of this counterexample.
\end{proof}

Thus the actual BT-generated formulation has existence and uniqueness
with the listed graded comparison, whereas the larger formulation
written on scan 023 is false.  Neither conclusion makes Newton slopes
a complete invariant of the sections: their exact classification remains
Theorem~\ref{cat:J1}, with the rank-one refinement in
Proposition~\ref{cat:J4}.

\subsection{The ICM programme: the precise historical comparison}
\begin{proposition}[Known, with the primary formulations compared]\label{cat:J9}
The ``mysterious functor'' lead is correct for reconstruction of
Barsotti--Tate isogeny data through crystalline representations.
It does not identify Fontaine's crystalline category with the literal
all-Hodge--Tate category of 161-5/023.
The relevant ICM article is \emph{Groupes de Barsotti-Tate et cristaux},
Nice 1970, tome~1, pp.~431--436 (published in 1971), not the separately
published Montreal course whose title adds \emph{de Dieudonn\'e}.
\end{proposition}
\begin{proof}[Numbered comparison of the opened statements]
\begin{enumerate}
\item Grothendieck's \S5, p.~435, Remarks~2--3
\cite{catGrothendieckICM1970}, ask for the Galois representation from
the filtered Dieudonne space and for the converse reconstruction and
essential image.  Fontaine's crystalline equivalence in
\cite[\S5.2]{catFontaine1982}, together with its BT compatibility
in \S6.2, provides these inverse constructions on the admissible
category.  One must choose consistently between covariant and
contravariant Dieudonne modules; the reconstruction assertion is
unchanged by the corresponding dualization.
\item The intrinsic description of that essential image requires a
further theorem.  Colmez--Fontaine's Theorem~A and \S4.1 identify
weak admissibility with admissibility
\cite{catColmezFontaine2000}.  Kisin's Corollary~(2.2.6)
\cite{catKisin2006} identifies the crystalline weight-$[0,1]$
representations with rational Tate modules of BT groups.  These
statements together answer the BT essential-image problem at the
isogeny level for the perfect-residue-field scope stated above.
Theorem~\ref{cat:J7} then supplies the precise uniqueness of the
tensor extension from actual BT objects with the specified graded
comparison.  It is not a classification by Newton slopes alone.
\item This connection is also documented by the authors of the modern
theorem: Colmez--Fontaine's Introduction, p.~3 of the opened author
version, identifies Grothendieck's questions on the ``foncteur
myst\'erieux'' as an origin of the theory, and its reference~[23] is
exactly the ICM article above.  Thus the historical link does not
depend on the later study's assertion alone.
\item The final paragraph of ICM p.~435 asks for analogous
comparisons in higher cohomology and over more general complete bases.
Those broader requests are not all the assertion of the 1982 BT
comparison theorem.  Likewise, the S3 obstruction in
Theorem~\ref{cat:J8} proves that the larger category literally written
in the dossier cannot simply be substituted into the modern theorem.
The defensible identification is of the reconstruction programme and
its crystalline solution, with these categorical hypotheses retained.
\end{enumerate}
\end{proof}

